\documentclass[12pt]{article}
\usepackage{geometry}
\geometry{margin=1.4in}
\usepackage{setspace}
\setstretch{1.4}
\usepackage{amsmath}
\usepackage{amssymb}

\title{The Boundary of Displacement\\
\large What This Framework Does NOT Explain}
\author{Diego Rincón}
\date{}

\begin{document}

\maketitle

\begin{abstract}
The displacement framework explains systems by their ground state $S_0$ (intended, designed, natural), their attractor $S^*$ (where they settle), and the cost to move between them. The adversarial critique is sharp: $S_0$ is chosen post-hoc, making the framework unfalsifiable. Whatever persists is claimed to hold $S_0$ at cost $C_{\text{maintenance}}$; whatever decays is claimed to drift toward $S^*$; whatever lies between is claimed to be displacement aging. The frame supplies its own coordinates retroactively, admitting no contrary observation. This paper grants the critique its full force and states exactly where it lands: the framework does NOT fit systems with no ground state to return to, systems where return is literally impossible, systems with emergent attractors that have no ancestral antecedent, systems that punish return attempts, systems governed by physics without agents, and systems locked by path-dependence. These boundaries are not defects to hide. They are the framework's real domain. The parts that survive this audit are the parts that do not depend on $S_0$ existing: the renewal-reward theorem (pure optimization independent of ground state), the empirical structure (displacement as observable, return cost as measurable), and the maintenance theorem's cube-root law (valid in any system with convex return cost). This paper is the one that saves the corpus from crankery by stating its limits.
\end{abstract}

\section{The Adversarial Attack}

The strongest critique of the displacement framework is this: \textit{$S_0$ is normative and retroactively defined. There is no way to falsify the model because the framework is unfalsifiable.}

Here is how it goes.

Suppose a system persists at some state $S^*$ for a long time. The framework says: ``This is displacement. The system is paying $C_{\text{maintenance}}$ to hold its ground state $S_0$ at cost, and has chosen to stay at $S^*$ rather than return.'' But how do we know $S_0$ is the right choice? Answer: because the system keeps paying $C_{\text{maintenance}}$. How do we know it's really paying $C_{\text{maintenance}}$ and not just stable? Answer: because it's displaced from $S_0$. The circle is complete. No observation can break it.

Suppose instead the system decays and collapses. The framework says: ``The displacement aged. The return cost scaled past any maintenance the system was prepared to pay, and $C_{\text{return}} > C_{\text{maintenance}}$, so the system irreversibly drifted to $S^*$.'' But how do we know $S_0$ was the ground state to begin with? Answer: because the system's collapse shows it was. The reasoning is tautological.

Or suppose the system occupies some intermediate state, neither fully $S_0$ nor fully $S^*$. The framework calls this ``displacement aging''—the system is partway from $S_0$ to $S^*$, held there by costs that make the mid-point locally stable. But what if the system is not displacing at all? What if this intermediate state \textit{is} the ground state, and the framework is simply mislabeling its coordinates retroactively?

The adversarial reviewer's conclusion: the framework is not falsifiable because $S_0$ is chosen post-hoc to match whatever the system is observed doing. The frame supplies its own coordinates. No observation contradicts it.

This critique is correct. It lands directly. The framework \textit{does} have this vulnerability. The job of this paper is not to deny it but to locate exactly where the vulnerability lives and where it does not.

\section{Where the Framework Does NOT Fit}

The framework fails in six classes of systems. Each failure points to a real boundary of applicability.

\subsection{1. Pure Extractors: Systems with No Prior Ground State}

A pure extractor is a system built from the start to extract value from something and offer nothing in return. Many colonial economies, industrial agriculture, and fossil fuel systems fit this description.

The framework assumes every system has a $S_0$---an intended, designed, or natural state it was meant to occupy. But consider a plantation economy built by slaveholders from day one. The ``ground state'' was never ``flourishing colonies with free people sharing wealth.'' The ground state was extraction. The entire system was architected for displacement from the outset. There is no $S_0$ to return to because there never was one.

In such a system, the framework's language breaks down. We cannot say the system is ``held at $S^*$ by paying $C_{\text{maintenance}}$'' because there is no implied contrast with a ground state. The system is functioning exactly as designed. Calling it a displacement is not wrong in spirit, but it requires importing a counterfactual ground state ($S_0$: what these lands and people should be if justly treated) that is external to the system's actual design. The framework works only when $S_0$ is immanent—internal to the system's logic—not when it is imposed from outside moral judgment.

This is a real boundary. For systems built on extraction from the start, the framework's explanatory power depends on injecting an external value judgment into the system's coordinates. That is philosophy, not description. The framework does not describe pure extractors; it judges them.

\subsection{2. Irreversible Thresholds: Systems Past the Point of No Return}

A system passes an irreversible threshold when $S_0$ no longer exists as a reachable state, not because the cost has scaled high, but because the prerequisites for $S_0$ have been destroyed.

Consider extinction. The framework says a species stays at a degraded attractor $S^*$ while $C_{\text{return}} > C_{\text{maintenance}}$. But if the species is extinct, $S_0$ is not expensive to return to. It is impossible. There is no cost curve that connects to it. The framework's entire geometry assumes a continuous cost landscape. Extinction eliminates the landscape. $C_{\text{return}} \to \infty$ is not the right description; $S_0$ is not on the map anymore.

The same applies to dead cultures, burned forests, and depleted aquifers. Once the ecological, cultural, or physical substrate is gone, return is not expensive—it is undefined. The framework's grammar does not fit.

More subtly, irreversible thresholds can be social. A form of knowledge can be forgotten so completely that re-learning it is not harder than learning it was; it is learning from scratch. A tradition can be broken across generations so that return would require not continuing a cost but inventing a new ground state. The framework assumes return means retracing a path. Some systems have no path left to retrace.

\subsection{3. Emergent Attractors: Systems Where $S^*$ Was Never $S_0$}

The framework assumes every attractor $S^*$ is a displacement from some prior ground state $S_0$. But some attractors emerge de novo. They were not displacements from anywhere.

Consider new technology. The ground state of human agency did not include the ability to compute at scale, communicate globally, or process data at machine speed. When new technologies emerge, they create new attractors that were never ground states. The human attention economy did not displace from some natural state into algorithmic curation; the algorithm is a new invention, and the attractor it creates was not a return to something humans ever did at scale before.

The same holds for new social norms. Democracy as currently practiced did not displace from some ancestral egalitarianism into representation; representation is a distinct invention with its own logic. To call it displacement is to invent a counterfactual past.

The framework is built for understanding how existing systems fall into wrong attractors. It is not built for understanding how new systems emerge. For emergent attractors, the framework's language of return is misleading. There is nothing to return to. The attractor is not a displacement; it is an invention.

\subsection{4. Adversarial Return: Systems That Punish the Attempt to Return}

The framework assumes return to $S_0$ costs $C_{\text{return}}(\tau)$ and is then available. But in some systems, the act of trying to return triggers defenses that escalate the cost.

Consider algorithmic feed systems. If you try to exit the algorithmic attractor (stop using the platform), the system's business model punishes you with loss of access to your network. If you try to reduce your algorithm exposure (switch to chronological feed), the platform's design makes it increasingly hard (chronological feeds are less visible, logged out, deprecated). The system is adversarial to your return attempt. The return cost is not a fixed function $C_{\text{return}}(\tau)$; it is a function that \textit{adapts as you try to return}. The attractor punishes attempts to leave.

The same applies to abusive relationships, addictive drugs, and certain kinds of social control. The return cost is not independent of the attempt to return. Trying to leave makes the cost higher. This is not just expensive return; it is \textit{adaptive} expensive return, where $S^*$ fights to keep you there.

The framework's model assumes passive cost landscape. It breaks when the landscape is an adversary.

\subsection{5. Systems Without Agents: Physics and Automatic Processes}

The framework's defining inequality is that a system stays at $S^*$ while ``$C_{\text{return}} > C_{\text{maintenance}}$.'' This language presupposes an agent who is choosing or being forced to stay. It presupposes someone is paying the cost.

But consider a falling stone. It is displaced from its ground state (at rest) into motion (its attractor). The framework would ask: is the stone paying $C_{\text{maintenance}}$ to stay in motion, or is it staying because return is expensive? The question is nonsensical. The stone is not paying anything. There is no agent. Gravity is not choosing between states; gravity is a law.

More generally, any pure physical process without agency—thermodynamic flows, weather systems, planetary orbits—fit the displacement model's geometry but not its agency logic. You can describe an avalanche as displaced from the mountain's $S_0$ (at rest) into motion toward $S^*$ (at the bottom). But the avalanche is not choosing to stay displaced. There is no cost being paid. The framework's explanatory power here is purely metaphorical.

This is a subtle boundary. The framework's mathematics can be applied to any system with a quasi-stable state and a cost function. But the framework's interpretation—the story it tells about agents and choices—requires that someone or something is actively paying to stay at $S^*$. Without agency, the story is empty.

\subsection{6. Hysteresis and Path-Dependence: Systems Locked by Infrastructure}

The framework distinguishes two cases: (a) systems that \textit{can} return to $S_0$ but choose not to because the cost is high, and (b) systems that \textit{cannot} return because the path has been blocked.

But real infrastructure locks are not cleanly in either category. Consider a transportation network built entirely around cars. The ground state was presumably walkable neighborhoods with mixed transit. The attractor is car-dependent sprawl. Can you return? In principle, yes—rebuild the neighborhoods, restore the transit. But the cost is not just high; it is compounded by hysteresis. Tearing out the car infrastructure creates a temporary state (rubble, disconnection) that is worse than the car-dependent attractor. You have to pass through intermediate states with even higher costs to reach $S_0$. The cost landscape has peaks that make the path impassable, not because $S_0$ is gone, but because the local topography blocks the route.

Similarly, social locks. Once a dominant group has consolidated power through extraction, the path to equality passes through a state (loss of power for the dominant group, temporary chaos) that makes the transition attractively worse than the stable extraction. The return cost includes not just the action but the temporary inversion cost.

The framework handles irreversibility as $C_{\text{return}}(\tau) = f + b\tau^{1+k}$ growing without bound. But it does not cleanly capture the local-peak problem: the cost rises not monotonically toward $S^*$ but with intermediate peaks that trap systems in locally stable but globally wrong states. Path-dependence is not just expensive return; it is return made topographically impossible by intermediate barriers.

\section{The Scope of Failure}

These six cases share a common structure: they all involve some deviation from the framework's core assumptions about $S_0$.

\begin{enumerate}
\item \textbf{Pure extractors} have no $S_0$ immanent in the system.
\item \textbf{Irreversible thresholds} have an $S_0$ that no longer exists.
\item \textbf{Emergent attractors} have no prior $S_0$ at all.
\item \textbf{Adversarial return} has a $S_0$ but an adaptive defense against reaching it.
\item \textbf{Systems without agents} have a $S_0$ but no one paying to stay away.
\item \textbf{Path-dependent locks} have a $S_0$ and an expensive return, but the path is blocked by intermediate peaks.
\end{enumerate}

Each of these is a real limit. The framework does not explain these systems. It can describe their geometry but not their dynamics or their narrative.

And the adversarial reviewer is right about the epistemological consequence: in these domains, the framework is vulnerable to the charge of unfalsifiability. If you can always redraw $S_0$ to match what the system is observed doing, then the framework is not testable. It is post-hoc coordinate assignment.

\section{What Survives the Audit}

But the framework does not stand or fall as a whole. Parts of it are independent of $S_0$ existing, and those parts survive the critique.

\subsection{The Renewal-Reward Theorem}

The foundation of the maintenance theorem is pure optimization theory. The governing equation is:
$$g(\tau)=\frac{a\,\delta^2\,\tau^2}{6}+\frac{f}{\tau}+b\,\tau^{k}.$$

This equation describes the long-run average cost per unit time for a ``reset every $\tau$'' policy in any system with:
\begin{itemize}
\item Convex aging cost $D(\xi) = \tfrac{1}{2}a\xi^2$
\item Drift rate $\delta$ toward some attractor
\item Fixed cost $f$ to reset
\item Super-linear irreversibility cost $b\tau^{1+k}$ in elapsed time
\end{itemize}

None of these require $S_0$ to be specially chosen or immanent. They apply to any system with these cost structures, whether or not we agree on what the ground state should be.

The cube-root law $\tau^* = \left(\frac{3f}{a\delta^2}\right)^{1/3}$ is pure economics. It is the economic-order-quantity rule applied to maintenance scheduling. The fact that it works regardless of what $S_0$ actually is means that the maintenance theorem's empirical content is decoupled from the ground-state problem.

If you want to \textit{know} whether a system is holding a ground state optimally, you do not need to agree on what the ground state is. You measure $f$, $\delta$, $a$, and observe $\tau$. If $\tau$ is near $\tau^*$, the system is optimization-aligned. If $\tau$ is far from $\tau^*$, the system is not. This is a testable claim.

The renewal-reward theorem survives because it is about cost structure, not ground state.

\subsection{The Empirical Structure}

The framework defines displacement as an observable quantity: $\xi$, distance from $S_0$, increases over time, is felt as cost $D(\xi)$, and can be reversed at cost $C_{\text{return}}(\tau)$.

This structure is independent of whether $S_0$ is chosen post-hoc. Here is why.

In any system you can identify:
\begin{itemize}
\item A state the system was in (or was designed for). Call it $S_0$.
\item A state it is actually in. Call it $S^*$.
\item A distance metric $\xi$ between them.
\item A felt cost of that distance: $D(\xi)$.
\item A cost to reduce that distance: $C_{\text{return}}(\tau)$.
\end{itemize}

Even if you later discover you chose $S_0$ wrong, these observables remain. The felt cost is real. The return cost is real. The distance is measurable. You have empirical content regardless of whether $S_0$ was correctly identified.

The critique that $S_0$ is retroactively defined does not invalidate the empirical measurements. It just means the interpretation of what $S_0$ means is contestable. The measurement is not.

This is important. It means the framework has an empirical core that does not depend on the philosophical correctness of ground state choice. You can test whether a system is paying convex aging costs, whether return is super-linear in time, whether the maintenance interval is optimal. These are testable. The ground state itself might be chosen wrong, but the cost structure around it is observable.

\subsection{The Maintenance Theorem's Cube-Root Law}

As stated above, the cube-root law is a pure optimization result. It holds in any system where the cost to hold a state against drift is convex, and the cost to reset is fixed plus super-linear in elapsed time.

This is true regardless of whether the system should be at $S_0$ or somewhere else.

The law says: if you want to minimize long-run average cost, given a fixed cost $f$ and a convex aging cost, reset at an interval proportional to $f^{1/3}$. The law does not care why you want to stay where you are. It cares only about the cost structure.

Applications:
\begin{itemize}
\item Optimal inventory is ordered on a schedule proportional to $f^{1/3}$ where $f$ is the order cost.
\item Optimal maintenance is performed on a schedule proportional to $f^{1/3}$ where $f$ is the setup cost of any restoration.
\item Optimal return (to any state, whether $S_0$ or otherwise) is scheduled proportional to $f^{1/3}$.
\end{itemize}

The theorem is about structure. It survives ground-state critique because it does not care what you are returning to. It only cares that you want to return, and at what cost.

\section{Three Concrete Tests for Falsifiability}

If the framework is to be made falsifiable beyond the parts that are simply optimization theory, it must make predictions that can be wrong. Here are three such tests.

\subsection{Test 1: Prospective Prediction in a Domain with No $S_0$}

Consider a system where the ground state is genuinely contested or absent: new social media platforms, nascent AI systems, or emerging technologies.

The test: make a prospective prediction about the system's behavior without choosing $S_0$ post-hoc. Specifically, predict the return cost as a function of time spent in the current attractor, based only on the cost structure's curvature and the assumed drift rate. If the actual return cost tracks the prediction, the framework has predictive power even without correct ground-state identification.

Example: A new attention-capture algorithm is deployed. If you model its drift rate and its convex aging cost, you can predict how hard it will be to ``escape'' after 6 months, 1 year, 2 years. If your prediction holds, the framework works. If the return cost is linear instead of convex, or sublinear, the framework fails.

This test decouples the framework from ground-state choice. It tests only whether the cost structure and return dynamics are as claimed.

\subsection{Test 2: Adversarial Return}

Identify a system where return attempts are plausibly punished adaptively (algorithmic platforms, abusive relationships, addiction systems).

The test: measure the cost to leave as a function of (a) time spent in the system and (b) the person's prior engagement level. If the system is adversarial to return, the cost should increase non-monotonically: it should rise sharply when you try to leave, beyond what simple irreversibility would predict.

Prediction: in non-adversarial systems, $C_{\text{return}}(\tau)$ is monotonic in $\tau$ alone. In adversarial systems, $C_{\text{return}}(\tau, I)$, where $I$ is prior engagement, should show $\frac{\partial^2 C}{\partial \tau \partial I} > 0$: the marginal cost of time spent rises with engagement. If this interaction term is absent, the adversarial mechanism is not active. If it is present, the framework's assumption of passive cost landscape is violated.

This test directly probes the critique that some systems fight return. If you can measure the adaptation, you can falsify the passive-landscape assumption.

\subsection{Test 3: Irreversibility and the Point of No Return}

In a system past an irreversible threshold (extinct species, dead cultures, depleted resources), measure $C_{\text{return}}$ explicitly and check whether it is infinite or merely very large.

Prediction: for systems with survivable ground state, $C_{\text{return}}(\tau)$ should be finite but convex-increasing. For systems past irreversible thresholds, $C_{\text{return}}(\tau)$ should be undefined or infinite (there is no path back), not just very expensive.

Test: take a system claimed to be at an irreversible threshold (fishery collapse, extinct language, genocide's cultural erasure). Estimate the cost to restore it to anything resembling $S_0$. If the cost is finite and super-linear, the threshold is not truly irreversible; the framework applies. If the cost is undefined (you cannot restore what does not exist in any living form), the framework does not apply.

This test directly probes the boundary between expensive return and impossible return.

\section{What This Means for the Corpus}

The falsifiability audit yields a clear result: the displacement framework explains \textit{bounded classes of systems well} and \textit{fails gracefully in others}.

It explains systems where:
\begin{itemize}
\item A ground state can be identified, whether immanent or external
\item Return is expensive but possible
\item Agents are actively paying costs to stay at attractors
\item The cost landscape is static (not adversarial)
\item Path-dependence is not locally prohibitive
\end{itemize}

It does not explain (or explains poorly) systems where:
\begin{itemize}
\item No ground state ever existed (pure extractors)
\item The ground state no longer exists (irreversible thresholds)
\item The attractor is emergent, not a displacement from a prior state
\item The cost landscape adapts to return attempts
\item No agents are active in the system
\item Intermediate barriers make the path impassable
\end{itemize}

This is not a defeat. This is a domain specification. Every theory has boundaries. The strength of a theory is knowing where they are.

The parts of the framework that survive the audit are the parts worth keeping:

1. The \textbf{renewal-reward theorem} is universally true for any system with convex aging cost and super-linear return cost. It is not falsified by the critique because it makes no claims about ground states.

2. The \textbf{empirical structure} (displacement as observable, return cost as measurable) is valid even if ground-state choice is contested. You can make measurements regardless.

3. The \textbf{maintenance theorem} is an optimization result. It holds wherever convex costs and scheduled returns are the decision problem, regardless of what you are returning to.

4. The \textbf{three levers} (frequency, fixed cost, drift) describe valid policy interventions in any system where those parameters are controllable.

5. The \textbf{irreversibility condition} ($C_{\text{return}} > C_{\text{maintenance}}$ traps systems at wrong attractors) is valid where the cost landscape is well-defined, which is not everywhere.

The corpus survives by narrowing its claims. It ceases to be a universal explanation of all human experience, and becomes what it can be: a framework for understanding how systems held in intentional states resist return, and what the cost structure of that resistance looks like.

That is still immense work.

\section{Conclusion: The Necessary Boundary}

The adversarial reviewer's critique is correct: $S_0$ is normative, retroactively defined, and vulnerable to unfalsifiability. The displacement framework does not have a principled way to choose ground states. It can describe their cost landscape only after they are chosen.

This is not hidden. It is the core of the framework: it asks ``given that a system is at $S^*$ and we believe it should be at $S_0$, what is the cost of return?'' The ``we believe'' is not a bug. It is a feature. The framework does not tell you what to value. It tells you what the cost is if you do value something and want to return to it.

But this means the framework has limits. It does not apply to systems where $S_0$ is absent or impossible. It does not apply to systems without agency. It does not apply to systems where the cost landscape is adversarial. It does not apply where the path is blocked by intermediate barriers. And it does not apply to emergent systems with no ancestral ground state.

These boundaries are not weaknesses to apologize for. They are the framework's true shape. A theory that knows its limits is stronger than one that pretends to none.

The framework asks one question: ``Given a state you value, how much does it cost to stay there or return to it?'' That is the right question for a bounded domain. It is not the question for systems where the valued state never existed, or has become impossible, or is under active attack. For those systems, other frameworks apply.

The corpus survives by accepting its boundaries.

And for the systems where the framework does apply—where ground state is identifiable, return is possible, agents are active, and the cost landscape is static—it explains the structure of displacement, maintenance, and return in a way no other framework does.

That is enough.

\end{document}
